\section*{Introduction}

The statistical toolkit of population genomics---nucleotide diversity, Fst, site frequency spectra, haplotype homozygosity, runs of homozygosity---has been refined over decades. Yet the biological questions these statistics are designed to answer increasingly demand analytical integration that the tools implementing them cannot provide. Has domestication relaxed purifying selection uniformly across all subpopulations of a crop species, and by how much? Is protein-coding divergence between ecotypes elevated relative to synonymous divergence, and where in the genome? Do multiple independent selection statistics converge on the same gene at a candidate locus? These questions are not answerable by any single existing tool---not because the underlying statistics are unavailable, but because current tools cannot condition statistics on functional annotations, cannot compose results from separate analyses, and discard computed values after each run.

These limitations are architectural, rooted in the flat-file, matrix-based paradigm shared by all widely used population genomics software.

\textbf{First, matrix-based computation is inherently wasteful.} Tools such as scikit-allel\cite{miles2024scikit} load entire genotype matrices ($N$ samples $\times$ $M$ variants) into memory, scaling as $O(N \times M)$. For 3{,}000 samples and 29~million variants, this implies $\sim$87~billion matrix entries---most read, decompressed, and discarded for each analysis. VCFtools\cite{danecek2011variant} re-parses the same VCF text for every statistic. PLINK~2\cite{chang2015second} mitigates this with binary formats but still couples each computation to file I/O. Every tool pays the full cost of data access for every analysis, regardless of whether the same data were accessed moments before.

\textbf{Second, computed results are ephemeral.} When scikit-allel computes iHS, the result is a NumPy array that vanishes when the session ends. VCFtools writes a text file; PLINK writes a \texttt{.fst} file. None of these results are queryable, filterable, or composable. Asking ``which variants show iHS~$> 2$ \emph{and} lie in a window with Fst~$> 0.3$?'' requires custom coordinate-matching scripts across output files. There is no persistent analytical record; the gap between computing a result and being able to use it in a follow-up question is as wide as computing it from scratch.

\textbf{Third, functional annotation is disconnected from statistical computation.} Connecting a selection signal to its target genes---the ``last mile''---requires exporting coordinates, running bedtools intersect\cite{quinlan2010bedtools}, matching against VEP\cite{mclaren2016ensembl} output, and manual curation. More fundamentally, no existing tool can compute a statistic \emph{restricted to variants of a given functional class}. The ratio $\pi_N/\pi_S$---a direct measure of the efficacy of purifying selection\cite{lu2006cost}---requires filtering a VCF by consequence type, computing diversity on each filtered subset, and repeating for every population and chromosome: at minimum four tool invocations and three intermediate files per combination. An annotation-conditioned genome scan (e.g., sliding-window Fst among only missense variants) is simply not possible with any current tool.

These are not merely inconveniences. They represent blind spots: entire classes of biological investigation that are theoretically straightforward but operationally inaccessible. The cost of domestication ($\pi_N/\pi_S$) has been studied in individual species-pair comparisons\cite{lu2006cost}, but a systematic genome-wide quantification across all subpopulations of a major crop---requiring hundreds of annotation-conditioned calls---has not been performed because no tool makes it practical. Similarly, whether protein-coding Fst exceeds synonymous Fst at specific population boundaries (evidence for adaptive protein divergence) has been examined in candidate regions but never via whole-genome annotation-conditioned scans.

Here we present GraphPop, an analytical engine comprising 12 stored procedures that run inside a Neo4j\cite{neo4j} graph database. GraphPop replaces the flat-file matrix paradigm with a graph-native architecture that addresses all three limitations and, critically, opens analytical pathways that lead to new biological discoveries:

\begin{enumerate}
\item \textbf{Graph-native computation.} Pre-aggregated allele-count arrays on Variant nodes yield $O(V \times K)$ complexity independent of sample count for six FAST~PATH procedures. Bit-packed genotypes (2~bits/sample) provide 87\% memory reduction for six FULL~PATH procedures. GraphPop maintains constant $\sim$150~MB memory versus 4.9~GB peak for scikit-allel, with speedups up to 1{,}614$\times$.

\item \textbf{Persistent, queryable analytical record.} Computed statistics---iHS scores, Fst values, ROH segments, window summaries---are written to graph nodes as permanent properties. Multi-statistic convergence (e.g., high XP-EHH \emph{and} PBS \emph{and} negative Fay \& Wu's $H$ at the same locus) is discoverable by a single query over already-computed values, with no re-computation and no file merging.

\item \textbf{Annotation-conditioned queries.} Every procedure accepts \texttt{consequence}, \texttt{pathway}, or gene filters resolved by graph edge traversal. Computing $\pi_N$, synonymous Fst, or a pathway-restricted genome scan is a one-parameter change---enabling systematic investigation of annotation-dependent evolutionary dynamics at genome scale.
\end{enumerate}

We validate GraphPop on the 1000 Genomes Project\cite{1000genomes2015global} (chr22: 3{,}202 samples, 1{,}070{,}401 variants; full genome: all autosomes) and then apply all 12 procedures to the 3K Rice Genomes Project\cite{wang2018genomic} (3{,}024 accessions, 29.6M SNPs, 12 subpopulations). The rice analysis demonstrates that graph-native investigation uncovers biology inaccessible to existing tools: annotation conditioning reveals, to our knowledge for the first time at this scale, that \emph{all 12 rice subpopulations} show genome-wide relaxed purifying selection ($\pi_N/\pi_S > 1.0$), and that missense Fst exceeds synonymous Fst at the temperate--tropical japonica boundary (ratio 1.11)---direct evidence of adaptive protein divergence during ecological specialization. The persistent analytical record enables convergent multi-statistic identification of known domestication genes (GW5, Hd1, PROG1, Wx, OsC1) via single queries. Pre-built Neo4j graph databases for both rice 3K and 1000 Genomes are publicly available for reuse. A companion platform, GraphMana, handles data management and export (separate publication), and an MCP server enables AI-agent-driven genomic investigation.
